\chapter*{Zusammenfassung}
\addcontentsline{toc}{chapter}{Zusammenfassung}
\selectlanguage{ngerman}

Zum 1.~Januar 2026 ist der CO\textsubscript{2}-Grenzausgleichsmechanismus
(CBAM) der Europäischen Union von einer reinen Berichtspflicht in die
definitive Phase übergegangen: Importeure müssen seither Zertifikate für die
in den erfassten Waren gebundenen Emissionen abgeben. Die vorliegende Arbeit
untersucht, ob dieser Übergang die europäischen Beschaffungsströme messbar
verschoben hat.

Die Identifikationsstrategie nutzt den Umstand, dass der Anwendungsbereich
über eine explizite Liste von Codes der Kombinierten Nomenklatur definiert
ist. Erfasste Warengruppen werden mit unmittelbar benachbarten, nicht
erfassten Codes verglichen --- ähnliche Güter, ähnliche Abnehmer, ähnliche
Konjunkturabhängigkeit, aber keine Zertifikatspflicht. Auf monatlichen
Comext-Einfuhrdaten von Eurostat wird ein
Differenz-von-Differenzen-Ansatz geschätzt. Betrachtet werden drei
Zielgrößen: Einfuhrmengen, die Umschichtung der Importanteile zwischen
Herkunftsländern unterschiedlicher Emissionsintensität sowie Durchschnittswerte
je Mengeneinheit als Näherung für die Preisüberwälzung.

[Ergebnisse werden nach Abschluss der Schätzung ergänzt.] Die Analyse ist so
angelegt, dass eine tatsächliche Emissionsverlagerung von blo\ss{}er
Umschichtung sauberer Produktionseinheiten unterschieden werden kann.

\medskip
\noindent\textbf{Schlagwörter:} \thesisKeywords

\selectlanguage{english}
